% -----------------------------------------------------------------------------
% Additional QED observables in strict SI units
% -----------------------------------------------------------------------------

轻费米子过程已经展示了交换、湮灭、阈值和相同粒子统计。现在换一个问题检验同一套 QED 振幅：如果散射的一端越来越重、越来越接近静态外源，相对论性单光子交换应怎样连续回到熟悉的 Coulomb 散射？这个极限同时检验外态归一化、重源展开和 Dirac 自旋修正，而不再引入新的相互作用。

以下仍从\eqnrefs{eq:2to2-cross-section-master-si-r9} 的二体相空间主式出发，并始终把 Mandelstam 变量视为物理四动量的平方；因此 $s,t,u$ 的量纲均为动量平方，截面的面积量纲由显式的 $\hbar^2$ 提供。

\begin{exbox}[Mott 与 Rutherford 极限]
考虑质量 $M_N$、电荷 $Ze$ 的重核，并在完整电子--原子核单光子交换之后取 $M_Nc\gg |\mathbf p|$。约化树级振幅为
\begin{equation}
 \mathcal M
 =\frac{Zg_{\rm em}^2}{t}
 [\bar u(p')\gamma^\mu u(p)]
 [\bar U(P')\gamma_\mu U(P)].
 \label{eq:mott-reduced-amplitude-si-r10}
\end{equation}
重核电流的主导项为
\begin{equation}
 \bar U(P')\gamma^\mu U(P)
 =2M_Nc\,\delta_0^\mu
 +\order\!\left(\frac{|\mathbf p|}{M_Nc}\right),
 \label{eq:heavy-nucleus-current-si-r10}
\end{equation}
因此核端只提供近静态时间分量，而电子端仍保留完整 Dirac 自旋结构。定义
\begin{equation*}
 \beta_{\rm rel}\equiv\frac{|\mathbf p|c}{E},
\end{equation*}
并在重核反冲的一阶修正可忽略时，得到 Mott 截面
\begin{equation}
 \boxed{
 \frac{\dd\sigma_{\rm Mott}}{\dd\Omega}
 =\frac{Z^2\hbar^2\alpha^2}
 {4|\mathbf p|^2\beta_{\rm rel}^2\sin^4(\theta/2)}
 \left[1-\beta_{\rm rel}^2\sin^2\frac\theta2\right].}
 \label{eq:mott-cross-section-si-r10}
\end{equation}
其中 $\hbar^2/|\mathbf p|^2$ 显式给出面积量纲。再取非相对论极限 $|\mathbf p|\simeq m_{\mathrm e} v$、$\beta_{\rm rel}\simeq v/c$，便得到 Rutherford 公式
\begin{equation}
 \boxed{
 \frac{\dd\sigma_{\rm Ruth}}{\dd\Omega}
 =\left[
 \frac{Ze^2}{16\pi\varepsilon_0E_{\rm kin}}
 \right]^2
 \frac1{\sin^4(\theta/2)}.}
 \label{eq:rutherford-SI-r10}
\end{equation}
因此 $\sin^{-4}(\theta/2)$ 的前向增强来自无质量光子的 $1/t$ 极点，而 Mott 因子 $1-\beta_{\rm rel}^2\sin^2(\theta/2)$ 来自电子 Dirac 电流的自旋与相对论结构。


\medskip

\eqnrefs{eq:mott-reduced-amplitude-si-r10} 已把电子--重核的一光子交换写成两条电流的乘积。先利用 $M_Nc\gg|\mathbf p|$ 把核电流约化为时间分量，可得到保留电子相对论自旋效应的 Mott 截面；随后再施加 $v/c\ll1$，这一结果必须连续退化为\eqnrefs{eq:rutherford-SI-r10}。这里补出这两个受控极限之间的自旋迹、重核展开和 SI 系数化简。



由\eqnrefs{eq:mott-reduced-amplitude-si-r10}，在 $M_Nc\gg|\mathbf p|$ 时使用
\begin{equation*}
 \bar U(P')\gamma^\mu U(P)
 =2M_Nc\,\delta_0^\mu
 +\order\!\left(\frac{|\mathbf p|}{M_Nc}\right).
\end{equation*}
电子端的未偏振时间分量为
\begin{align*}
 N_e
 &\equiv\frac12\sum_{s,s'}
 |\bar u_{s'}(p')\gamma^0u_s(p)|^2\\
 &=\frac12\Tr[(\slashed p'+m_{\mathrm e} c)\gamma^0
 (\slashed p+m_{\mathrm e} c)\gamma^0]\\
 &=2\left[2p'^0p^0-(p'\!\cdot p-m_{\mathrm e}^2c^2)\right].
\end{align*}
弹性条件 $p'^0=p^0=E/c$ 以及
\begin{equation*}
 p'\!\cdot p=\frac{E^2}{c^2}-|\mathbf p|^2\cos\theta,
 \qquad
 \frac{E^2}{c^2}=|\mathbf p|^2+m_{\mathrm e}^2c^2
\end{equation*}
给出
\begin{equation*}
 N_e
 =\frac{4E^2}{c^2}
 \left[1-\beta_{\rm rel}^2\sin^2\frac\theta2\right].
\end{equation*}
因此
\begin{equation*}
 \overline{|\mathcal M|^2}
 =\frac{16Z^2g_{\rm em}^4M_N^2E^2}{t^2}
 \left[1-\beta_{\rm rel}^2\sin^2\frac\theta2\right]
 +\order\!\left(\frac{|\mathbf p|}{M_Nc}\right).
\end{equation*}
重核极限还有
\begin{equation*}
 s=M_N^2c^2\left[1+\order\!\left(\frac{E}{M_Nc^2}\right)\right],
 \qquad
 \frac{|\mathbf p_f|}{|\mathbf p_i|}=1+\order\!\left(\frac{E}{M_Nc^2}\right),
\end{equation*}
故二体截面主式化为
\begin{equation*}
 \frac{\dd\sigma}{\dd\Omega}
 =\frac{Z^2\hbar^2g_{\rm em}^4E^2}
 {4\pi^2c^2t^2}
 \left[1-\beta_{\rm rel}^2\sin^2\frac\theta2\right].
\end{equation*}
弹性运动学给 $t=-4|\mathbf p|^2\sin^2(\theta/2)$；再代入 $g_{\rm em}^2=4\pi\alpha$ 与 $\beta_{\rm rel}=|\mathbf p|c/E$，即得\eqnrefs{eq:mott-cross-section-si-r10}。

最后取 $|\mathbf p|\simeq m_{\mathrm e} v$、$\beta_{\rm rel}\simeq v/c$：
\begin{equation*}
 \frac{\dd\sigma}{\dd\Omega}
 \longrightarrow
 \frac{Z^2\alpha^2\hbar^2c^2}
 {4m_{\mathrm e}^2v^4\sin^4(\theta/2)}.
\end{equation*}
由
\begin{equation*}
 \alpha=\frac{e^2}{4\pi\varepsilon_0\hbar c},
 \qquad
 E_{\rm kin}=\frac12m_{\mathrm e} v^2,
\end{equation*}
其系数逐步化为
\begin{align*}
 \frac{Z^2\alpha^2\hbar^2c^2}{4m_{\mathrm e}^2v^4}
 &=\frac{Z^2}{4m_{\mathrm e}^2v^4}
 \left(\frac{e^2}{4\pi\varepsilon_0}\right)^2\\
 &=\left[
 \frac{Ze^2}{16\pi\varepsilon_0E_{\rm kin}}
 \right]^2,
\end{align*}
从而得到\eqnrefs{eq:rutherford-SI-r10}。
\end{exbox}

\begin{figure}[htbp]
 \centering
 \includegraphics[width=.74\linewidth]{mott_rutherford_angular.png}
 \caption{Mott/Rutherford 角分布的相对论修正。各曲线在 $\theta=30^\circ$ 处归一化；$\beta_{\rm rel}$ 增大时，Dirac 电流因子 $1-\beta_{\rm rel}^2\sin^2(\theta/2)$ 对大角散射的抑制逐渐明显，而小角增强仍由 $\sin^{-4}(\theta/2)$ 主导。}
 \label{fig:mott-rutherford-angular-r6}
\end{figure}

\begin{note}[前向奇异性的物理截断]
\eqnrefs{eq:mott-cross-section-si-r10,eq:rutherford-SI-r10} 假定无限程、未屏蔽的 Coulomb 势，平面波渐近态以及理想角分辨。真实原子核受到原子电子屏蔽；有限束斑和有限相互作用时间给出最小可分辨动量转移；有限探测器接受角也会把前向锥积分成有限值。对晶体、二维材料或纳米结构，更一般的处理不是人为给 $q$ 加截断，而是使用结构化介电响应或推迟 Green 张量，使屏蔽和边界条件进入传播核本身。
\end{note}

同一个国际单位制约化 QED 顶角 $-\ii g_{\rm em}\gamma^\mu$ 在不同外态指定下产生不同通道。$s$ 通道的虚四动量是入射总四动量，$t$ 通道是四动量转移，$u$ 通道则对应交换一个末态标签后的四动量转移。交叉对称性的实质，是同一个解析 Green 函数在不同外部四动量上作入射/出射解析延拓；不能把图形的几何旋转替代自旋子、统计符号和物理区域运动学的解析延拓。

